\documentclass[11pt]{article}
\usepackage[a4paper,margin=0.8in]{geometry}
\usepackage{calc}
\usepackage{array,longtable,booktabs}
\usepackage{xcolor}
\usepackage{hyperref}
\usepackage{enumitem}
\usepackage{fancyvrb}
\usepackage{fvextra}
\usepackage{framed}
\usepackage{color}
\fvset{breaklines=true, breakanywhere=true}
\DefineVerbatimEnvironment{Highlighting}{Verbatim}{breaklines,breakanywhere,commandchars=\\\{\}}
\usepackage{color}
\usepackage{fancyvrb}
\newcommand{\VerbBar}{|}
\newcommand{\VERB}{\Verb[commandchars=\\\{\}]}
\DefineVerbatimEnvironment{Highlighting}{Verbatim}{commandchars=\\\{\}}
% Add ',fontsize=\small' for more characters per line
\usepackage{framed}
\definecolor{shadecolor}{RGB}{248,248,248}
\newenvironment{Shaded}{\begin{snugshade}}{\end{snugshade}}
\newcommand{\AlertTok}[1]{\textcolor[rgb]{0.94,0.16,0.16}{#1}}
\newcommand{\AnnotationTok}[1]{\textcolor[rgb]{0.56,0.35,0.01}{\textbf{\textit{#1}}}}
\newcommand{\AttributeTok}[1]{\textcolor[rgb]{0.13,0.29,0.53}{#1}}
\newcommand{\BaseNTok}[1]{\textcolor[rgb]{0.00,0.00,0.81}{#1}}
\newcommand{\BuiltInTok}[1]{#1}
\newcommand{\CharTok}[1]{\textcolor[rgb]{0.31,0.60,0.02}{#1}}
\newcommand{\CommentTok}[1]{\textcolor[rgb]{0.56,0.35,0.01}{\textit{#1}}}
\newcommand{\CommentVarTok}[1]{\textcolor[rgb]{0.56,0.35,0.01}{\textbf{\textit{#1}}}}
\newcommand{\ConstantTok}[1]{\textcolor[rgb]{0.56,0.35,0.01}{#1}}
\newcommand{\ControlFlowTok}[1]{\textcolor[rgb]{0.13,0.29,0.53}{\textbf{#1}}}
\newcommand{\DataTypeTok}[1]{\textcolor[rgb]{0.13,0.29,0.53}{#1}}
\newcommand{\DecValTok}[1]{\textcolor[rgb]{0.00,0.00,0.81}{#1}}
\newcommand{\DocumentationTok}[1]{\textcolor[rgb]{0.56,0.35,0.01}{\textbf{\textit{#1}}}}
\newcommand{\ErrorTok}[1]{\textcolor[rgb]{0.64,0.00,0.00}{\textbf{#1}}}
\newcommand{\ExtensionTok}[1]{#1}
\newcommand{\FloatTok}[1]{\textcolor[rgb]{0.00,0.00,0.81}{#1}}
\newcommand{\FunctionTok}[1]{\textcolor[rgb]{0.13,0.29,0.53}{\textbf{#1}}}
\newcommand{\ImportTok}[1]{#1}
\newcommand{\InformationTok}[1]{\textcolor[rgb]{0.56,0.35,0.01}{\textbf{\textit{#1}}}}
\newcommand{\KeywordTok}[1]{\textcolor[rgb]{0.13,0.29,0.53}{\textbf{#1}}}
\newcommand{\NormalTok}[1]{#1}
\newcommand{\OperatorTok}[1]{\textcolor[rgb]{0.81,0.36,0.00}{\textbf{#1}}}
\newcommand{\OtherTok}[1]{\textcolor[rgb]{0.56,0.35,0.01}{#1}}
\newcommand{\PreprocessorTok}[1]{\textcolor[rgb]{0.56,0.35,0.01}{\textit{#1}}}
\newcommand{\RegionMarkerTok}[1]{#1}
\newcommand{\SpecialCharTok}[1]{\textcolor[rgb]{0.81,0.36,0.00}{\textbf{#1}}}
\newcommand{\SpecialStringTok}[1]{\textcolor[rgb]{0.31,0.60,0.02}{#1}}
\newcommand{\StringTok}[1]{\textcolor[rgb]{0.31,0.60,0.02}{#1}}
\newcommand{\VariableTok}[1]{\textcolor[rgb]{0.00,0.00,0.00}{#1}}
\newcommand{\VerbatimStringTok}[1]{\textcolor[rgb]{0.31,0.60,0.02}{#1}}
\newcommand{\WarningTok}[1]{\textcolor[rgb]{0.56,0.35,0.01}{\textbf{\textit{#1}}}}
\setlist[itemize]{leftmargin=*}
\setlist[enumerate]{leftmargin=*}
\hypersetup{colorlinks=true,linkcolor=blue,urlcolor=blue}
\providecommand{\tightlist}{%
  \setlength{\itemsep}{0pt}\setlength{\parskip}{0pt}}
\title{R Workshop -- Day 4 -- Probability, Sampling, Simulation and Inference}
\author{Applied R for Statistical Teaching, Research and Reproducible Analysis}
\date{}
\begin{document}
\maketitle
\tableofcontents
\newpage
\hypertarget{r-workshop-day-4-probability-sampling-simulation-and-inference}{%
\section{R Workshop -- Day 4 -- Probability, Sampling, Simulation and
Inference}\label{r-workshop-day-4-probability-sampling-simulation-and-inference}}

\textbf{Session 4 of 7 \textbar{} Duration: 3 hours} \textbf{Programme:}
Applied R for Statistical Teaching, Research and Reproducible Analysis

\hypertarget{learning-objectives}{%
\subsection{Learning objectives}\label{learning-objectives}}

By the end of this session you will be able to:

\begin{enumerate}
\def\labelenumi{\arabic{enumi}.}
\tightlist
\item
  Use the \texttt{d}/\texttt{p}/\texttt{q}/\texttt{r}
  distribution-function pattern for at least the normal and binomial
  distributions.
\item
  Set a random seed and explain why it matters for reproducible
  simulation.
\item
  Run a simulation experiment and demonstrate the Central Limit Theorem
  with repeated sampling.
\item
  Estimate and interpret a confidence interval in plain language.
\item
  Choose an appropriate test based on study design and run it on the
  anchor dataset.
\item
  Report statistical and practical significance without overstating
  certainty.
\end{enumerate}

\hypertarget{capstone-link}{%
\subsection{Capstone link}\label{capstone-link}}

Today produces one inferential answer to your capstone research
question: typically, whether post-training scores improved relative to
pre-training scores, and/or whether scores differ between two groups
(e.g.~completers versus non-completers).

\hypertarget{segment-plan-3-hours}{%
\subsection{Segment plan (3 hours)}\label{segment-plan-3-hours}}

\begin{longtable}[]{@{}
  >{\raggedright\arraybackslash}p{(\columnwidth - 4\tabcolsep) * \real{0.3000}}
  >{\raggedleft\arraybackslash}p{(\columnwidth - 4\tabcolsep) * \real{0.4000}}
  >{\raggedright\arraybackslash}p{(\columnwidth - 4\tabcolsep) * \real{0.3000}}@{}}
\toprule\noalign{}
\begin{minipage}[b]{\linewidth}\raggedright
Segment
\end{minipage} & \begin{minipage}[b]{\linewidth}\raggedleft
Time
\end{minipage} & \begin{minipage}[b]{\linewidth}\raggedright
Purpose
\end{minipage} \\
\midrule\noalign{}
\endhead
\bottomrule\noalign{}
\endlastfoot
Opening and recap & 10 min & Recap Day 3 graphics, preview today's
inferential question \\
Concept bridge & 25 min & Probability functions, sampling variation, the
CLT \\
Guided coding & 60 min & Simulation, confidence interval, one hypothesis
test \\
Participant practice & 45 min & Run inference on the anchor dataset \\
Interpretation and reporting & 25 min & Plain-language statistical
reporting \\
Evidence log and capstone update & 15 min & Save the inference
notebook/script \\
\end{longtable}

\begin{center}\rule{0.5\linewidth}{0.5pt}\end{center}

\hypertarget{part-1-concept-bridge-core-workshop-activity}{%
\subsection{Part 1 -- Concept bridge (Core workshop
activity)}\label{part-1-concept-bridge-core-workshop-activity}}

\hypertarget{the-dpqr-pattern}{%
\subsubsection{\texorpdfstring{The
\texttt{d}/\texttt{p}/\texttt{q}/\texttt{r}
pattern}{The d/p/q/r pattern}}\label{the-dpqr-pattern}}

Every distribution in R follows a naming convention, illustrated for the
normal distribution:

\begin{longtable}[]{@{}
  >{\raggedright\arraybackslash}p{(\columnwidth - 4\tabcolsep) * \real{0.3333}}
  >{\raggedright\arraybackslash}p{(\columnwidth - 4\tabcolsep) * \real{0.3333}}
  >{\raggedright\arraybackslash}p{(\columnwidth - 4\tabcolsep) * \real{0.3333}}@{}}
\toprule\noalign{}
\begin{minipage}[b]{\linewidth}\raggedright
Prefix
\end{minipage} & \begin{minipage}[b]{\linewidth}\raggedright
Meaning
\end{minipage} & \begin{minipage}[b]{\linewidth}\raggedright
Example
\end{minipage} \\
\midrule\noalign{}
\endhead
\bottomrule\noalign{}
\endlastfoot
\texttt{d} & density (height of the curve at x) &
\texttt{dnorm(70,\ mean\ =\ 65,\ sd\ =\ 10)} \\
\texttt{p} & cumulative probability, P(X \textless= x) &
\texttt{pnorm(70,\ mean\ =\ 65,\ sd\ =\ 10)} \\
\texttt{q} & quantile, the x for a given cumulative probability &
\texttt{qnorm(0.95,\ mean\ =\ 65,\ sd\ =\ 10)} \\
\texttt{r} & random draws from the distribution &
\texttt{rnorm(10,\ mean\ =\ 65,\ sd\ =\ 10)} \\
\end{longtable}

The same four letters apply to \texttt{binom} (binomial), \texttt{unif}
(uniform), \texttt{t} (t-distribution), and others.

\hypertarget{why-set.seed}{%
\subsubsection{\texorpdfstring{Why
\texttt{set.seed()}}{Why set.seed()}}\label{why-set.seed}}

\begin{Shaded}
\begin{Highlighting}[]
\FunctionTok{set.seed}\NormalTok{(}\DecValTok{2024}\NormalTok{)}
\FunctionTok{rnorm}\NormalTok{(}\DecValTok{3}\NormalTok{)}
\end{Highlighting}
\end{Shaded}

Running the line above twice with the same seed gives the same three
numbers. Without \texttt{set.seed()}, every run of a simulation differs,
which makes results impossible to check or reproduce. Set the seed once,
near the top of the script, before the first random draw.

\hypertarget{sampling-variation-and-the-central-limit-theorem-clt}{%
\subsubsection{Sampling variation and the Central Limit Theorem
(CLT)}\label{sampling-variation-and-the-central-limit-theorem-clt}}

A single sample mean is one estimate of the population mean; a different
sample would give a different estimate. The CLT says that, as sample
size grows, the distribution of the \emph{sample mean} across repeated
samples becomes approximately normal, regardless of the shape of the
original population -- which is why so much of classical inference
relies on the normal distribution.

\begin{center}\rule{0.5\linewidth}{0.5pt}\end{center}

\hypertarget{part-2-guided-coding-core-workshop-activity}{%
\subsection{Part 2 -- Guided coding (Core workshop
activity)}\label{part-2-guided-coding-core-workshop-activity}}

\begin{Shaded}
\begin{Highlighting}[]
\FunctionTok{library}\NormalTok{(tidyverse)}

\NormalTok{lecturers }\OtherTok{\textless{}{-}} \FunctionTok{read\_csv}\NormalTok{(}\StringTok{"Data/processed/anchor\_dataset\_clean.csv"}\NormalTok{, }\AttributeTok{show\_col\_types =} \ConstantTok{FALSE}\NormalTok{)}
\FunctionTok{set.seed}\NormalTok{(}\DecValTok{2024}\NormalTok{)}
\end{Highlighting}
\end{Shaded}

\hypertarget{one-simulation-experiment}{%
\subsubsection{2.1 One simulation
experiment}\label{one-simulation-experiment}}

\begin{Shaded}
\begin{Highlighting}[]
\CommentTok{\# Simulate 1,000 lecturers\textquotesingle{} post{-}training scores under an assumed}
\CommentTok{\# population mean of 65 and sd of 12, purely to see typical sample noise.}
\NormalTok{simulated\_scores }\OtherTok{\textless{}{-}} \FunctionTok{rnorm}\NormalTok{(}\DecValTok{1000}\NormalTok{, }\AttributeTok{mean =} \DecValTok{65}\NormalTok{, }\AttributeTok{sd =} \DecValTok{12}\NormalTok{)}
\FunctionTok{mean}\NormalTok{(simulated\_scores)}
\FunctionTok{sd}\NormalTok{(simulated\_scores)}
\FunctionTok{hist}\NormalTok{(simulated\_scores, }\AttributeTok{main =} \StringTok{"Simulated post{-}training scores"}\NormalTok{, }\AttributeTok{xlab =} \StringTok{"Score"}\NormalTok{)}
\end{Highlighting}
\end{Shaded}

\hypertarget{demonstrating-the-clt-with-repeated-samples}{%
\subsubsection{2.2 Demonstrating the CLT with repeated
samples}\label{demonstrating-the-clt-with-repeated-samples}}

\begin{Shaded}
\begin{Highlighting}[]
\NormalTok{sample\_means }\OtherTok{\textless{}{-}} \FunctionTok{replicate}\NormalTok{(}\DecValTok{2000}\NormalTok{, \{}
\NormalTok{  one\_sample }\OtherTok{\textless{}{-}} \FunctionTok{sample}\NormalTok{(lecturers}\SpecialCharTok{$}\NormalTok{post\_score[}\SpecialCharTok{!}\FunctionTok{is.na}\NormalTok{(lecturers}\SpecialCharTok{$}\NormalTok{post\_score)], }\AttributeTok{size =} \DecValTok{30}\NormalTok{)}
  \FunctionTok{mean}\NormalTok{(one\_sample)}
\NormalTok{\})}

\FunctionTok{hist}\NormalTok{(sample\_means, }\AttributeTok{main =} \StringTok{"Distribution of sample means (n = 30, 2000 repeats)"}\NormalTok{,}
     \AttributeTok{xlab =} \StringTok{"Sample mean post{-}training score"}\NormalTok{)}
\FunctionTok{sd}\NormalTok{(sample\_means) }\CommentTok{\# the standard error, empirically}
\end{Highlighting}
\end{Shaded}

Compare \texttt{sd(sample\_means)} to
\texttt{sd(lecturers\$post\_score,\ na.rm\ =\ TRUE)\ /\ sqrt(30)} -- the
two should be close, which is the CLT's standard-error formula made
concrete.

\hypertarget{one-confidence-interval}{%
\subsubsection{2.3 One confidence
interval}\label{one-confidence-interval}}

\begin{Shaded}
\begin{Highlighting}[]
\NormalTok{post\_scores }\OtherTok{\textless{}{-}}\NormalTok{ lecturers}\SpecialCharTok{$}\NormalTok{post\_score[}\SpecialCharTok{!}\FunctionTok{is.na}\NormalTok{(lecturers}\SpecialCharTok{$}\NormalTok{post\_score)]}
\NormalTok{n }\OtherTok{\textless{}{-}} \FunctionTok{length}\NormalTok{(post\_scores)}
\NormalTok{sample\_mean }\OtherTok{\textless{}{-}} \FunctionTok{mean}\NormalTok{(post\_scores)}
\NormalTok{sample\_se }\OtherTok{\textless{}{-}} \FunctionTok{sd}\NormalTok{(post\_scores) }\SpecialCharTok{/} \FunctionTok{sqrt}\NormalTok{(n)}

\NormalTok{ci\_lower }\OtherTok{\textless{}{-}}\NormalTok{ sample\_mean }\SpecialCharTok{{-}} \FunctionTok{qt}\NormalTok{(}\FloatTok{0.975}\NormalTok{, }\AttributeTok{df =}\NormalTok{ n }\SpecialCharTok{{-}} \DecValTok{1}\NormalTok{) }\SpecialCharTok{*}\NormalTok{ sample\_se}
\NormalTok{ci\_upper }\OtherTok{\textless{}{-}}\NormalTok{ sample\_mean }\SpecialCharTok{+} \FunctionTok{qt}\NormalTok{(}\FloatTok{0.975}\NormalTok{, }\AttributeTok{df =}\NormalTok{ n }\SpecialCharTok{{-}} \DecValTok{1}\NormalTok{) }\SpecialCharTok{*}\NormalTok{ sample\_se}
\FunctionTok{c}\NormalTok{(ci\_lower, ci\_upper)}

\CommentTok{\# Or, directly from a t{-}test:}
\FunctionTok{t.test}\NormalTok{(post\_scores)}\SpecialCharTok{$}\NormalTok{conf.int}
\end{Highlighting}
\end{Shaded}

\textbf{Plain-language template:} ``We are 95\% confident that the true
mean post-training score for the population this sample represents lies
between \texttt{ci\_lower} and \texttt{ci\_upper}.''

\hypertarget{selecting-and-running-one-test}{%
\subsubsection{2.4 Selecting and running one
test}\label{selecting-and-running-one-test}}

Test choice follows study design, not habit:

\begin{longtable}[]{@{}
  >{\raggedright\arraybackslash}p{(\columnwidth - 2\tabcolsep) * \real{0.5000}}
  >{\raggedright\arraybackslash}p{(\columnwidth - 2\tabcolsep) * \real{0.5000}}@{}}
\toprule\noalign{}
\begin{minipage}[b]{\linewidth}\raggedright
Design
\end{minipage} & \begin{minipage}[b]{\linewidth}\raggedright
Typical test
\end{minipage} \\
\midrule\noalign{}
\endhead
\bottomrule\noalign{}
\endlastfoot
One numeric variable vs.~a fixed value & One-sample t-test \\
Same individuals measured twice (pre/post) & Paired t-test \\
Two independent groups, one numeric outcome & Independent (two-sample)
t-test \\
Two categorical variables & Chi-square test of independence \\
\end{longtable}

For the capstone question -- did scores improve from pre- to
post-training -- the design is paired:

\begin{Shaded}
\begin{Highlighting}[]
\NormalTok{paired\_data }\OtherTok{\textless{}{-}}\NormalTok{ lecturers }\SpecialCharTok{\%\textgreater{}\%} \FunctionTok{filter}\NormalTok{(}\SpecialCharTok{!}\FunctionTok{is.na}\NormalTok{(pre\_score), }\SpecialCharTok{!}\FunctionTok{is.na}\NormalTok{(post\_score))}

\FunctionTok{t.test}\NormalTok{(paired\_data}\SpecialCharTok{$}\NormalTok{post\_score, paired\_data}\SpecialCharTok{$}\NormalTok{pre\_score, }\AttributeTok{paired =} \ConstantTok{TRUE}\NormalTok{)}
\end{Highlighting}
\end{Shaded}

\hypertarget{reporting-statistical-and-practical-significance}{%
\subsubsection{2.5 Reporting statistical and practical
significance}\label{reporting-statistical-and-practical-significance}}

\begin{Shaded}
\begin{Highlighting}[]
\NormalTok{paired\_result }\OtherTok{\textless{}{-}} \FunctionTok{t.test}\NormalTok{(paired\_data}\SpecialCharTok{$}\NormalTok{post\_score, paired\_data}\SpecialCharTok{$}\NormalTok{pre\_score, }\AttributeTok{paired =} \ConstantTok{TRUE}\NormalTok{)}
\NormalTok{mean\_gain }\OtherTok{\textless{}{-}} \FunctionTok{mean}\NormalTok{(paired\_data}\SpecialCharTok{$}\NormalTok{post\_score }\SpecialCharTok{{-}}\NormalTok{ paired\_data}\SpecialCharTok{$}\NormalTok{pre\_score)}
\end{Highlighting}
\end{Shaded}

\textbf{Reporting template:} ``Post-training scores were on average
\texttt{mean\_gain} points higher than pre-training scores (paired
t-test, t = \ldots, df = \ldots, p = \ldots). This difference is
{[}statistically significant / not statistically significant{]} at the
5\% level, and represents a {[}meaningful / modest{]} practical gain on
a 0-100 scale.''

Statistical significance (a small p-value) and practical significance (a
gain large enough to matter to a lecturer or programme manager) are
different claims -- report both.

\begin{center}\rule{0.5\linewidth}{0.5pt}\end{center}

\hypertarget{part-3-participant-practice-core-workshop-activity}{%
\subsection{Part 3 -- Participant practice (Core workshop
activity)}\label{part-3-participant-practice-core-workshop-activity}}

\begin{enumerate}
\def\labelenumi{\arabic{enumi}.}
\tightlist
\item
  Compute a 95\% confidence interval for mean attendance.
\item
  Run the paired t-test for pre- versus post-training scores and write
  the plain-language sentence above with your actual numbers filled in.
\item
  Choose one additional comparison relevant to your research question
  (e.g.~post-training score for completers versus non-completers, an
  independent two-sample t-test) and run it.
\item
  For your chosen test, check at least one assumption informally (e.g.~a
  histogram or \texttt{qqnorm()} of the relevant variable) and note any
  concern in a comment.
\item
  Save your code and your written interpretation as
  \texttt{Outputs/day4\_inference\_notebook.R} (or an
  \texttt{.Rmd}/\texttt{.qmd} if you prefer).
\end{enumerate}

A facilitator-reviewed solution is in
\texttt{Solution-Scripts/day4\_inference\_solution.R}.

\begin{center}\rule{0.5\linewidth}{0.5pt}\end{center}

\hypertarget{optional-demonstration}{%
\subsection{Optional demonstration}\label{optional-demonstration}}

\begin{itemize}
\tightlist
\item
  \textbf{Coverage simulation} -- repeat the confidence-interval
  calculation across many simulated samples and check what proportion
  actually contain the true mean (should be close to 95\%).
\item
  \textbf{Chi-square and Fisher's exact test} -- e.g.~is
  \texttt{completed} associated with \texttt{training\_track}?
\item
  \textbf{Non-parametric alternatives} -- Wilcoxon signed-rank test as a
  paired-t-test alternative when normality is doubtful.
\item
  \textbf{Effect-size calculations} -- Cohen's d via the
  \texttt{effectsize} package.
\item
  \textbf{Additional CLT demonstrations} -- repeat Part 2.2 with a
  clearly skewed variable (e.g.~\texttt{years\_teaching}) to show the
  CLT still holds for the \emph{sample mean} even when the raw variable
  is not normal.
\end{itemize}

\begin{Shaded}
\begin{Highlighting}[]
\CommentTok{\# Optional: chi{-}square test of independence}
\FunctionTok{chisq.test}\NormalTok{(}\FunctionTok{table}\NormalTok{(lecturers}\SpecialCharTok{$}\NormalTok{completed, lecturers}\SpecialCharTok{$}\NormalTok{training\_track))}
\end{Highlighting}
\end{Shaded}

\begin{center}\rule{0.5\linewidth}{0.5pt}\end{center}

\hypertarget{reference-or-take-home-material}{%
\subsection{Reference or take-home
material}\label{reference-or-take-home-material}}

\begin{itemize}
\tightlist
\item
  Assumption checklists for the t-test family (independence, approximate
  normality of the relevant quantity, for two-sample tests: variance
  assumptions and Welch's correction).
\item
  A short template for reporting a p-value and confidence interval
  consistently across the rest of the workshop.
\end{itemize}

\begin{center}\rule{0.5\linewidth}{0.5pt}\end{center}

\hypertarget{daily-deliverable}{%
\subsection{Daily deliverable}\label{daily-deliverable}}

A short inference notebook or script
(\texttt{Outputs/day4\_inference\_notebook.R} or equivalent) that
answers one inferential question linked to the capstone research
question, with a plain-language interpretation.

\hypertarget{evidence-log-entry}{%
\subsubsection{Evidence log entry}\label{evidence-log-entry}}

Complete the Day 4 row in \texttt{Workshop-Evidence-Log.md}: state your
inferential result and whether it supports, contradicts, or is
inconclusive about your initial impression from the Day 3 graphics.

\begin{center}\rule{0.5\linewidth}{0.5pt}\end{center}

\hypertarget{facilitator-notes}{%
\subsection{Facilitator notes}\label{facilitator-notes}}

\begin{itemize}
\tightlist
\item
  The paired t-test on \texttt{pre\_score}/\texttt{post\_score} should
  show a positive, statistically significant gain by construction of the
  synthetic dataset -- use this as your answer key for spot-checking
  participants' work.
\item
  Watch for participants computing a confidence interval or t-test on a
  column that still contains \texttt{NA} without filtering first;
  \texttt{t.test()} will error or silently drop rows depending on the
  call, so this is worth flagging explicitly.
\item
  If a participant chooses chi-square as their ``one additional
  comparison,'' that is acceptable even though it is listed as optional
  demonstration -- the boundary exists for time management, not as a
  hard rule.
\end{itemize}
\end{document}
